\section{Repeated labels and the outer-color dichotomy}
\label{sec:dichotomy}

We show that every cyclic progression either receives many occurrences
of each outer bit or admits one affine lift in the centered second
coordinates. The issue is repeated labels: a closest return to a heavy
label identifies a rational period and a small drift. Short periods
disappear under the dilation, while longer periods supply the eligible
patterns tested in \Cref{outer:coloring}.

\begin{lemma}[Affinity inside a small box]\label{dichotomy:affinity}
Let $z_j=a+jb$ in $\T^e$, for $0\le j<k$, and let
$J\subseteq\{0,\ldots,k-1\}$. Suppose real representatives
$Z_j\in\R^e$, $j\in J$, lie in one box whose coordinate intervals
have length at most $2H$. If $2kH<1$, there are $A,B\in\R^e$
such that $Z_j=A+jB$ for every $j\in J$.
\end{lemma}

\begin{proof}
For any $j_1<j_2<j_3$ in $J$, affinity modulo one implies that
\[
 (j_3-j_2)(Z_{j_2}-Z_{j_1})
 -(j_2-j_1)(Z_{j_3}-Z_{j_2})\in\Z^e.
\]
In each coordinate its absolute value is at most
$2H(j_3-j_1)\le2kH<1$, so it is zero. Thus every triple has the
same affine interpolation. If $|J|\ge2$, interpolate between its
smallest and largest indices to obtain $A,B$; the triple identity
places every remaining point on that interpolation. For at most two
indices the conclusion follows directly, including the empty case.
\end{proof}

Fix a cyclic progression $n_j=n_0+jd$, $0\le j<k$, with a heavy
label. In the next three lemmas the terms \emph{light} and \emph{heavy}
refer to this fixed progression.

\begin{lemma}[Closest return and drift]\label{dichotomy:drift}
Choose a heavy label and let $J$ be its set of occurrence indices,
with $m_*=|J|>k/M$. Let $h$ be the smallest positive difference
between two members of $J$, attained at $j'$ and $j'+h$.
Set
\[
 u=\hat x(n_{j'+h})-\hat x(n_{j'}),\qquad
 v=\tilde y(n_{j'+h})-\tilde y(n_{j'}).
\]
Then, for sufficiently large $k$,
\begin{equation}\label{dichotomy:drift-bounds}
 \begin{gathered}
 1\le h\le2M,\\
 |u_i|\le\frac{H_x}{m_*-1}\le\frac{2MH_x}{k},\qquad
 |v_i|\le\frac{2H}{m_*-1}\le\frac{4MH}{k}
 \quad(1\le i\le D),
 \end{gathered}
\end{equation}
and $v=\lambda u$ as real vectors.
\end{lemma}

\begin{proof}
For large $k$, the inequality $m_*>k/M$ gives
$m_*-1\ge k/(2M)$. Write $j_-=\min J$ and $j_+=\max J$.
The smallest of the $m_*-1$ consecutive gaps in $J$ satisfies
\[
 h\le\frac{j_+-j_-}{m_*-1}\le\frac{k-1}{m_*-1}\le2M.
\]
All vectors $(\hat x(n_j),\tilde y(n_j))$, $j\in J$, lie in
one full-label box, whose coordinate widths are at most $2H$ by
\Cref{mesh:width}. By \Cref{dichotomy:affinity}, their differences
are proportional to differences of indices. In particular the
difference at the chosen gap $h$ is $h/(j_+-j_-)$ times the difference
between the two extreme visits. The latter has first-coordinate
absolute differences at most $H_x$ and second-coordinate absolute
differences at most $2H$. This proves \Cref{dichotomy:drift-bounds}.

The identity $y=\lambda x$ on the torus implies
$v-\lambda u\in\Z^D$. Since $\lambda H_x=H$,
\[
 |v_i-\lambda u_i|\le\frac{6MH}{k}<1
\]
for large $k$. Every coordinate of this integer vector is therefore
zero.
\end{proof}

\begin{lemma}[Rational path and separated residues]\label{dichotomy:path}
With $h,u,v$ as in \Cref{dichotomy:drift}, there is
$t\in\{0,\ldots,h-1\}^D$ such that
\begin{equation}\label{dichotomy:path-equations}
 \begin{aligned}
 x(n_j)&=x(n_0)+\frac jh(t+u),\\
 y(n_j)&=y(n_0)+\frac jh(\lambda t+v)
 \end{aligned}
 \qquad\text{in }\T^D\quad(0\le j<k).
\end{equation}
Moreover $\gcd(t_1,\ldots,t_D,h)=1$, and indices incongruent
modulo $h$ have distinct full labels. Consequently, all occurrences
of any one full label lie in a single residue class modulo $h$.
\end{lemma}

\begin{proof}
Since $h x(d)=u$ modulo $\Z^D$, the vector
$t'=h\hat x(d)-u$ is integral. Reduce it coordinatewise modulo
$h$ to obtain $t$. The first equality in
\Cref{dichotomy:path-equations} follows; multiplication by $\lambda$
and the real identity $v=\lambda u$ give the second.

Suppose an integer $e>1$ divides $h$ and every $t_i$. It also divides
every $t'_i$, because $t'-t\in h\Z^D$. Hence the torus
displacements from index $j'$ to index $j'+h/e$ are $u/e$ and
$v/e$ in the first and second systems. The vector
\[
 (\hat x(n_{j'}),\tilde y(n_{j'}))+\frac1e(u,v)
\]
is a convex combination of the representatives at $j'$ and $j'+h$.
Their common full-label box is a product of half-open real intervals
and is convex. This vector therefore lies in that same box and is
the chosen representative of the sample at $j'+h/e$.
It gives an additional visit to the label at a strictly smaller
positive gap $h/e$, a contradiction. This proves the gcd assertion.

The torus vector $t/h$ has order exactly $h$: if $h$ divides all
$a t_i$, the gcd assertion and B\'ezout's identity imply $h\mid a$.
Take any $j_1,j_2\in\{0,\ldots,k-1\}$ incongruent modulo $h$.
In at least one coordinate the rational displacement
$(j_1-j_2)t_i/h$ has circle distance at least
$1/h$ from zero. The drift in that coordinate has absolute value at
most
\[
 \frac{k|u_i|}{h}\le\frac{2MH_x}{h}\le\frac{2MH}{h}.
\]
By the triangle inequality on the circle the actual displacement has
distance at least
\[
 \frac{1-2MH}{h}\ge\frac{1-2MH}{2M}>H\ge H_x
\]
for sufficiently large $k$. Two samples in the same $U$ interval
have circle distance at most $H_x$, so these samples cannot have
the same full label.
\end{proof}

\begin{lemma}[Drift relative to every heavy interval]
\label{dichotomy:heavy-width}
For every heavy index $j_*$ of the fixed progression and each
$1\le i\le D$,
\begin{equation}\label{dichotomy:relative-width}
 |v_i|\le\frac{2M}{k}\,|V(y_i(n_{j_*}))|.
\end{equation}
The heavy label here need not be the one used to choose $h$.
\end{lemma}

\begin{proof}
Let the full label at $j_*$ have multiplicity $m'>k/M$, and let
$j_-$ and $j_+$ be its extreme occurrence indices. By
\Cref{dichotomy:path}, $j_+-j_-=a h$ for an integer
$m'-1\le a<k/h$. Set $L_i=|V(y_i(n_{j_*}))|$. The real
representative difference
\[
 \Delta_i=\tilde y_i(n_{j_+})-\tilde y_i(n_{j_-})
\]
has $|\Delta_i|\le L_i\le2H$. From
\Cref{dichotomy:path-equations}, it is congruent to $a v_i$
modulo one. On the other hand,
\[
 |a v_i|\le\frac{4MH}{h},\qquad
 |\Delta_i-a v_i|\le2H+\frac{4MH}{h}<1
\]
for large $k$. The integer $\Delta_i-a v_i$ is zero. Therefore
\[
 |v_i|\le\frac{L_i}{a}\le\frac{L_i}{m'-1}
 \le\frac{2M}{k}L_i,
\]
as claimed. This argument uses only the two endpoint representatives;
it requires no assertion about the intervening path across a cut.
\end{proof}

\begin{theorem}[Outer-color dichotomy]\label{dichotomy:main}
For every sufficiently large $k$, the coloring $c_0$ from
\Cref{outer:pullback} has the following property. Every cyclic
progression $n_j=n_0+jd$, $0\le j<k$, with $d\ne0$ in $G$,
satisfies at least one of these alternatives:
\begin{enumerate}[label=\textup{(\roman*)}]
 \item each bit occurs among $c_0(n_0),\ldots,c_0(n_{k-1})$
 at least $\gamma k$ times;
 \item there are $A,B\in\R^D$ such that
 $\tilde y(n_j)=A+jB$ for all $0\le j<k$.
\end{enumerate}
The threshold on $k$ is absolute.
\end{theorem}

\begin{proof}
If the progression has at least $\delta k$ light positions,
\Cref{outer:coloring}\textup{(i)} supplies at least
$\delta k/4\ge\gamma k$ occurrences of each bit. We may therefore
assume that fewer than $\delta k$ positions are light. There is a
heavy label, and \Cref{dichotomy:drift,dichotomy:path,dichotomy:heavy-width}
provide $h,t,u,v$ with all the properties proved above.

\paragraph{Short periods.}
Suppose $h\le h_0$. The definition of the dilation gives
$h\mid\lambda$. Thus the second path in
\Cref{dichotomy:path-equations} has real step $w=v/h$ modulo one,
because $\lambda t/h$ is integral. Its total coordinate motion satisfies
\begin{equation}\label{dichotomy:short-motion}
 k|w_i|\le\frac{4MH}{h}<\frac12
\end{equation}
for large $k$. The congruence $w_i=y_i(d)$ modulo one also gives
$w_i\in q^{-i}\Z$: a representative of $y_i(d)$ is an integral
multiple of $q^{-i}$, and adding an integer preserves this lattice.
In particular,
\begin{equation}\label{dichotomy:step-lattice}
 w_i\ne0\quad\Longrightarrow\quad
 |w_i|\ge q^{-i}\ge q^{-D}\ge\alpha.
\end{equation}

Consider the real segment
$\tilde y_i(n_0)+\tau w_i$, $0\le\tau\le k-1$.
If the sampled centered representatives fail to equal
$\tilde y(n_0)+jw$, some coordinate has $w_i\ne0$ and this closed
segment meets a cut point in $1/2+\Z$. The word \emph{closed}
includes a cut met at the initial or final sample, so both directions
of motion and exact cut hits are covered. Each sampled point is
then within real distance $k|w_i|$ of such a cut point, giving
$\rho(y_i(n_j))\le k|w_i|$ for every $j$.
By \Cref{mesh:width,dichotomy:step-lattice},
\[
 |V(y_i(n_j))|
 \le2H\bigl(\rho(y_i(n_j))+\alpha\bigr)
 \le2H(k+1)|w_i|<|w_i|.
\]
Yet \Cref{dichotomy:short-motion} implies that any two distinct
sampled points in this coordinate have circle distance
$|j-j'|\,|w_i|\ge|w_i|$: the entire real displacement is less
than one half. They cannot belong to a common $V$ interval of width
less than $|w_i|$. All full labels would therefore be distinct,
contradicting the heavy label. Hence every sampled representative
equals $\tilde y(n_0)+jw$, proving alternative~\textup{(ii)}.

\paragraph{Longer periods.}
Suppose $h_0<h\le2M$. Partition the first
$h\floor{k/h}$ indices into consecutive full blocks of length $h$;
the remaining block has fewer than $h$ indices. A full block starting
at $j_0$ has exactly the form in \Cref{count:eligible}, with starting
points $x(n_{j_0}),y(n_{j_0})$ and the same $h,t,u,v$.
We verify its regularity and, when it contains a heavy index,
every condition of eligibility.

For a rotating coordinate $i$, let
$s_i=h/\gcd(h,\lambda t_i)>h_0$. Without the drift, the $h$
values
\[
 y_i(n_{j_0})+\frac{z\lambda t_i}{h},\qquad0\le z<h,
\]
run through an equally spaced grid of $s_i$ circle points, each
exactly $h/s_i$ times. The actual value at $z$ is within
$|v_i|\le4MH/k<\eta$ of this grid value. If the actual position is
nonregular in this coordinate, its grid value is within $2\eta$
of the cut. An arc of length $4\eta$ contains at most
$4\eta s_i+1$ points of an $s_i$-point grid. This bound remains
valid when endpoints are included. Consequently the fraction of
positions excluded by this coordinate is at most
$4\eta+1/s_i$. A union bound over rotating coordinates shows that
the fraction of nonregular positions is at most
\begin{equation}\label{dichotomy:regular-fraction}
 D\left(4\eta+\frac1{h_0}\right)
 =\frac4{1000}+\frac1D<\frac12.
\end{equation}
Thus each full block has at least $h/2$ regular positions. All labels
within the block are distinct by \Cref{dichotomy:path}, since its
indices have distinct residues modulo $h$. Also
$|u_i|\le H_x$ and $|v_i|\le H$ for all $i$ by
\Cref{dichotomy:drift-bounds}, for sufficiently large $k$.

Now use any full block containing a heavy index $j_*$.
Fix a stationary coordinate $i$ and write
$I_*=V(y_i(n_{j_*}))$, $L_*=|I_*|$. Its rational step
$\lambda t_i/h$ is integral, so for every index $j$ in the block,
\[
 \dist_{\T}\bigl(y_i(n_j),y_i(n_{j_*})\bigr)
 \le\frac{|j-j_*|}{h}|v_i|
 \le|v_i|\le\frac{2M}{k}L_*.
\]
For large $k$ this places $y_i(n_j)$ in the circular
$L_*$-neighborhood of $\overline{I_*}$. Applying
\Cref{mesh:neighbors} with $I=I_*$ and
$J=V(y_i(n_j))$ gives
\[
 |V(y_i(n_j))|\ge\frac{L_*}{16}\ge|v_i|,
\]
where the final inequality uses $2M/k\le1/16$ and
\Cref{dichotomy:relative-width}. This holds at every block position,
in particular at each regular position as eligibility requires.
The heavy index $j_*$ need not itself be regular: it serves only to
provide the interval $I_*$ and the relative drift estimate.
The neighborhood statement uses closures, so a cut crossing or an
endpoint equality causes no change in this verification.

Every full block containing a heavy index is therefore eligible.
By \Cref{outer:coloring}\textup{(ii)} and
\Cref{dichotomy:regular-fraction}, it contributes at least $h/8$
occurrences of each bit. A full block not used in this way consists
entirely of light positions. Since fewer than $\delta k$ positions
are light in the entire progression, and fewer than $h$ positions
lie in the remainder, the used full blocks cover at least
\[
 k-h-\delta k\ge\frac45 k
\]
positions for sufficiently large $k$, using $h\le2M=o(k)$.
Each bit consequently occurs at least $k/10\ge\gamma k$ times.
This proves alternative~\textup{(i)} and completes the dichotomy.
All inequalities imposed only absolute lower bounds on $k$.
\end{proof}
